%% kapa: exam=1 round=22 year=2011 period=2 form=A subject=appraisal_law subject_ko=부동산관계법규 seq=5 number=45
%% source: Q-Net 감정평가사 기출문제 (cst00309) artlSeq=5036790 file=extracted/1008/5036790/2011년 제22회 감정평가사 1차시험 기출문제/2교시.hwp
%% conversion: kordoc 4.24.51 -> tools/gen1.py (auto); review=proofread-gpt-astra; flags=none
\begin{question}{45}
국토의 계획 및 이용에 관한 법령상 지형도면에 관한 설명으로 옳지 않은 것은?
\begin{choices}
  \item 군수는 지형도면을 작성하면 도지사의 승인을 받아야 한다.
  \item 국토해양부장관이나 도지사는 도시관리계획을 직접 입안한 경우 관계 특별시장ㆍ광역시장, 시장 또는 군수의 의견을 들어 직접 지형도면을 작성할 수 있다.
  \item 자연녹지지역안의 지목이 대(垈)인 지역에 대해 축척 3천분의 1의 지형도를 사용하여 도시관리계획 결정을 고시한 경우에는 지형도면을 따로 작성할 필요가 없다.
  \item 지형도면의 고시를 해야함에도 불구하고 도시관리계획 결정의 고시일로부터 2년이 되는 날까지 지형도면의 고시가 없는 경우에는 그 2년이 되는 날의 다음 날에 그 도시관리계획 결정은 효력을 잃는다.
  \item 판례에 따르면, 도시계획결정고시의 도면만으로는 도시계획결정의 구체적인 범위나 개별토지의 도시계획선을 특정할 수 없으므로 도시계획결정 효력의 구체적ㆍ개별적인 범위는 지형도면고시에 의하여 확정된다.
\end{choices}
\end{question}
\answer{3}
